\section{Data}
\label{section:data}

GAIA-2 is trained on a large-scale internal dataset specifically curated to support the diverse demands of video generation for autonomous driving. The dataset comprises approximately 25 million video sequences, each spanning 2 seconds, collected between 2019 and 2024. Recordings were obtained across three countries—the United Kingdom, the United States, and Germany—to ensure coverage of geographically and environmentally diverse driving conditions.

To capture the complexity of real-world autonomous driving, data collection involved multiple vehicle platforms, including three different car models and two van types. Each vehicle was outfitted with either five or six cameras, configured to provide comprehensive 360-degree surround-view coverage. The camera systems varied in their capture frequencies—20 Hz, 25 Hz, and 30 Hz—introducing a range of temporal resolutions. This variability reflects the heterogeneous nature of sensor setups deployed in actual autonomous vehicles and supports GAIA-2’s ability to generalize across different input rates and hardware specifications.

An important characteristic of the dataset is the variability in camera placement throughout the data collection period. Over time, the positions and calibrations of the cameras were adjusted across platforms, introducing a broad spectrum of spatial configurations. This diversity provides a strong training signal for generalizing across different camera rigs, a key requirement for scalable synthetic data generation in the autonomous driving domain.

The dataset also encompasses a wide range of driving scenarios, including diverse weather conditions, times of day, road types, and traffic environments. To ensure coverage across this complexity, we explicitly balance the training data not just on individual features, but on their joint probability distribution. This approach enables a more representative learning signal by modeling realistic co-occurrences—e.g., specific lighting and weather conditions in certain geographies or behaviors unique to particular road types. To prevent redundant training samples, we enforce a minimum temporal stride between selected sequences, reducing the risk of duplication while maintaining a natural distribution.

For evaluation, we implement a geographically held-out validation strategy. Specific validation geofences are defined to exclude certain regions from the training set entirely. This ensures that model evaluation is performed on unseen locations, providing a more rigorous assessment of generalization performance across different environments.

In summary, this dataset provides a robust foundation for training GAIA-2. Its extensive temporal and spatial coverage, diversity in vehicle and camera configurations, and principled validation setup make it uniquely well-suited for developing generative world models capable of producing realistic and controllable driving video across varied real-world conditions.